\section{Central integrals marked at an arbitrary actual return}
\label{sec:marked-return-band}
A weight at a terminal or intermediate return cannot simply be smoothed
as a function of the initial state: composition with a long first-return
iterate need not preserve a uniform variation bound. Instead we change
the origin of the orbit to the marked return. The record then has a
backward and a forward part, while the weight is a function of the new
origin. Both parts can be estimated with the same forward collision
operator. This argument does not require invariance of the section under
time reversal, or mixing of the induced map.

All measures and normalizations in this section are as follows. The
measure $\nu$ is the probability on the full collision cylinder;
$\nu(Y_R^*)=c_*$, and $\nu_R^*=c_*^{-1}\nu|_{Y_R^*}$. For a complex
function $a$ that is zero almost everywhere outside $Y_R^*$ put
\begin{equation}\label{eq:marked-norms}
 M_a=\|a\|_\infty,\qquad V_a=\|a\|_{\mathrm{BV}},\qquad
 \alpha_a=\int_M a\dd\nu.
\end{equation}
For complex functions the variation norm is the sum of the norms of the
real and imaginary parts. It includes their $L^1$ norms and is taken
in the original collision coordinates, including the section boundary.
No upper bound for $V_a$ independent of $n$ is imposed below.

For $0\le k\le n$ define the finite complex measure and its centered,
rescaled transform by
\begin{align}
 \mu^{[k],a}_{n,R}(B)
   &=\int_{Y_R^*}a((F_R^*)^kx)\one_{\{J_{n,R}(x)\in B\}}\dd\nu(x),
       \label{eq:marked-measure}\\
 C^{[k],a}_{n,R}(v)
   &=\int_{Y_R^*}a((F_R^*)^kx)
          e^{iv\cdot(J_{n,R}(x)-n\bar G_R)/\sqrt n}\dd\nu(x).
       \label{eq:marked-transform}
\end{align}
Its mass is $\alpha_a$, which may be complex. The endpoints $k=0,n$
are the initial and terminal insertions. Taking $a=s_R$ gives the
original probability for every $k$. The physical frequency is always
$\omega=v/\sqrt n$.

\subsection{Changing the origin to the mark}
Let $N^-_{j,R}(y)$ be the collision distance to the $j$th preceding
section visit, so that
\[
 (F_R^*)^{-j}y=T_R^{-N^-_{j,R}(y)}y,\qquad N^-_{0,R}=0.
\]
The forward distance is $N^+_{j,R}=N_{j,R}$, with $N^+_{0,R}=0$.
These distances are finite almost everywhere, by recurrence and
invertibility. For nonnegative integers $r,s$ write
\begin{equation}\label{eq:two-sided-collision-sum}
 \mathsf H_{r,s,R}(y)=\sum_{j=-r}^{s-1}h_R(T_R^jy),
\end{equation}
where a sum over an empty interval is zero.

\begin{lemma}[Exact compensation about a marked return]
\label{lem:marked-compensation}
For almost every $y\in Y_R^*$,
\begin{equation}\label{eq:marked-compensation}
 J_{n,R}((F_R^*)^{-k}y)-n\bar G_R
   =\mathsf H_{N^-_{k,R}(y),N^+_{n-k,R}(y),R}(y).
\end{equation}
Consequently
\begin{equation}\label{eq:marked-recentering}
 C^{[k],a}_{n,R}(v)=\int_{Y_R^*}a(y)
 e^{iv\cdot\mathsf H_{N^-_{k,R}(y),N^+_{n-k,R}(y),R}(y)/\sqrt n}
 \dd\nu(y).
\end{equation}
\end{lemma}
\begin{proof}
The collision interval from $(F_R^*)^{-k}y$ up to, but excluding,
$(F_R^*)^{n-k}y$ has left endpoint $-N^-_{k,R}(y)$ and right endpoint
$N^+_{n-k,R}(y)$ in coordinates based at $y$. The sum of $f_R$ on this
interval is precisely the original four-coordinate record. Its negative
part contains exactly $k$ section visits, including the left endpoint
and excluding zero. Its nonnegative part contains exactly $n-k$ visits,
including zero if that part is nonempty and excluding the right endpoint.
Thus the sum of $\eta_R$ over the entire interval is $n$. Substitute
$h_R=f_R-\bar G_R\eta_R$. This proves the first identity, including
$k=0,n$. The actual first-return map preserves $\nu|_{Y_R^*}$, since it
preserves its normalization $\nu_R^*$. The change of variable
$y=(F_R^*)^kx$ gives the second identity. Only additivity and invariance
are used; the two parts are not asserted independent.
\end{proof}

\begin{lemma}[Forward and backward return-clock windows]
\label{lem:two-sided-clock-window}
For $j\ge0$, $b\ge2$, and $m_j=\lfloor j/c_*\rfloor$, uniformly in $R$,
\begin{equation}\label{eq:two-sided-clock-window}
 \nu_R^*\{|N^\pm_{j,R}-m_j|>b\}\le C\frac{j+b}{b^2}.
\end{equation}
Moreover the collision maximal estimate of
Lemma~\ref{lem:dyadic-collision-maximal} holds for an interval beginning
at any integer collision time, positive or negative, and in either order.
\end{lemma}
\begin{proof}
The case $j=0$ is immediate. For either sign let
$H_L^\pm=\sum_{i=1}^L\eta_R\circ T_R^{\pm i}$. Its mean under $\nu$
is $c_*L$. The variance is at most $CL$ by
Proposition~\ref{prop:bv-correlation}. For the negative sign the scalar
lag correlation is the same as for the positive sign by invariance.
Domination $\nu_R^*\le c_*^{-1}\nu$ gives
\[
 \int|H_L^\pm-c_*L|^2\dd\nu_R^*\le CL.
\]
The equivalences $N_j^\pm>L\iff H_L^\pm<j$ and
$N_j^\pm\le L\iff H_L^\pm\ge j$ hold for the actual visits. Take
integer thresholds within a fixed distance of $m_j+b$ and $m_j-b$.
When the latter threshold is negative, its lower-tail event is empty.
Otherwise the required deviation from $c_*L$ is at least $c_*b-O(1)$,
while $L\le C(j+b)$. Chebyshev's inequality proves the estimate for
$b$ above a fixed constant. For the remaining bounded range enlarge $C$
so that its right-hand side is at least one. In particular no assumption
$b\le j$ is needed.

For the maximal assertion, translate any deterministic integer interval
to a nonnegative one using invariance of $\nu$. Every subinterval still
has second moment bounded by its length, from the absolutely summable
collision correlations. The same dyadic prefix proof applies in the
reverse order. Domination gives its version under $\nu_R^*$ as well.
The assertion concerns fixed intervals, not a conditional distribution
of the past given the future.
\end{proof}

\begin{proposition}[Uniform stopping error about the mark]
\label{prop:marked-stopping}
For all $0\le k\le n$, $n\ge2$, and $v\in\R^4$,
\begin{align}
 &\left|C^{[k],a}_{n,R}(v)
  -\int_M a(y)e^{iv\cdot\mathsf H_{m_k,m_{n-k},R}(y)/\sqrt n}\dd\nu(y)
       \right|\notag\\
 &\hspace{12mm}\le C M_a\left(n^{-1/5}
             +|v|n^{-1/5}\log(2+n)\right).
       \label{eq:marked-stopping-error}
\end{align}
\end{proposition}
\begin{proof}
Use \eqref{eq:marked-recentering} and let
$b_n=\lceil n^{3/5}\rceil$. By Lemma~\ref{lem:two-sided-clock-window},
the union of the two events on which a clock differs from its
corresponding $m_j$ by more than $b_n$ has probability at most
$Cn^{-1/5}$ under $\nu_R^*$. Its contribution to the difference of
characteristic functions is at most $C M_an^{-1/5}$, since
$|a|\nu\le M_ac_*\nu_R^*$ on the section.

On the complement, the difference of the two collision sums is bounded
by the sum of two maxima on deterministic windows of length at most
$2b_n+2$ around their respective endpoints. The preceding maximal
estimate bounds its squared integral under $\nu$ by
$C b_n[\log(2+b_n)]^2$. It remains valid when a window crosses collision
time zero. Cauchy--Schwarz and $|e^{ix}-e^{iy}|\le|x-y|$ give
$C M_a|v|\sqrt{b_n/n}\log(2+b_n)$, which is the second term in
\eqref{eq:marked-stopping-error}. Nothing is estimated in an unbounded
stopped sum on the exceptional event. Finitely many initial values of
$n$ are included by increasing the constant.
\end{proof}

\subsection{A collision multiplier between two spectral powers}
Set $a_\delta=\mathsf S_\delta a$ and $h_{R,\delta}=\mathsf S_\delta h_R$.
The former need not be supported in the section. This is harmless:
\begin{equation}\label{eq:marked-smoothing-norms}
 \|a-a_\delta\|_1\le C\delta V_a,\qquad
 \|a_\delta\|_\infty\le M_a,\qquad
 \|M_{a_\delta}\|_{\mathcal B_j\to\mathcal B_j}
       \le C M_a\delta^{-2}.
\end{equation}
The last bound is Lemma~\ref{lem:density-norm-details}. In particular
its constant depends on the supremum norm, not on $V_a$.

\begin{lemma}[A marked deterministic collision integral]
\label{lem:marked-spectral}
Fix $c_0,C_0>0$. Suppose $c_0n\le r+s\le C_0n$ with $r,s\ge0$,
and let $q\ge1$ be an integer satisfying $2q\le r+s$. There are
$c,C>0$ and $\rho<1$ such that, whenever
\[
 |v|/\sqrt n\le c\delta^2,\qquad
 \delta^{-6}|v|^3/\sqrt n\le c,
\]
the following bound holds, with $\ell_\delta=1+|\log\delta|$:
\begin{align}
 &\left|\int_M a\,e^{iv\cdot\mathsf H_{r,s,R}/\sqrt n}\dd\nu
      -\alpha_a e^{-\frac{r+s}{2n}v^{\mathsf T}\Gamma_Rv}\right|
       \notag\\
 &\quad\le C\delta V_a+C M_a\left\{
 |v|\sqrt{\delta\ell_\delta}+|v|^2\delta\ell_\delta
 +\frac{|v|\delta^{-4}+|v|^3\delta^{-6}}{\sqrt n}
 +\delta^{-2}\rho^q+\frac{|v|q}{\sqrt n}
 +\frac{|v|^2q}{n}\right\}.
       \label{eq:marked-spectral-bound}
\end{align}
All constants are independent of the location of the mark and of $a$.
\end{lemma}
\begin{proof}
First replace $a$ by $a_\delta$. The cost is at most $C\delta V_a$.
Let $u_\delta=h_R-h_{R,\delta}$. Stationarity and
\eqref{eq:small-bv-variance} give
\[
 \int\left|\sum_{i=-r}^{s-1}u_\delta(T_R^iy)\right|^2\dd\nu(y)
       \le C(r+s)\delta\ell_\delta.
\]
The bounded weight $|a_\delta|\le M_a$ and Cauchy--Schwarz therefore
bound the cost of replacing $h_R$ by $h_{R,\delta}$ by
$C M_a|v|\sqrt{\delta\ell_\delta}$.

Put $z=v/\sqrt n$ and $L_z=\mathcal L_{R,\delta,z}$, using one of
the fixed local collision spaces. Its mass functional will be denoted
by $\ell$ to avoid confusion with block indices. Invariance and the
transfer convention give the exact chronological pairing
\begin{equation}\label{eq:marked-chronological-pairing}
 \int_M a_\delta(y)
 e^{iz\cdot\sum_{i=-r}^{s-1}h_{R,\delta}(T_R^iy)}\dd\nu(y)
       =\ell\bigl(L_z^s M_{a_\delta} L_z^r\nu\bigr).
\end{equation}
Indeed, starting at $x=T_R^{-r}y$, the rightmost power records collisions
$0,\ldots,r-1$, the multiplier is evaluated at $T_R^rx=y$, and the
leftmost power records collisions $r,\ldots,r+s-1$. Thus all weights
are evaluated at the prescribed states.

Suppose first that $r,s\ge q$. Write
$L_z^m=\lambda(z)^m\Pi(z)+E_m(z)$ with
$\|E_m(z)\|\le C\rho^m$. The hypotheses and
\eqref{eq:smooth-cubic-expansion} bound each real-frequency principal
power of length at most $r+s$ by a uniform constant. Expanding
\eqref{eq:marked-chronological-pairing}, every complementary term is
bounded by $C M_a\delta^{-2}(\rho^r+\rho^s)$; the term with both
complementary powers is included in this bound. For the principal
amplitude use $\Pi(0)=\nu\ell$ and \eqref{eq:marked-smoothing-norms}:
\begin{align*}
 \ell(\Pi(0)M_{a_\delta}\Pi(0)\nu)&=\alpha_a,\\
 |\ell(\Pi(z)M_{a_\delta}\Pi(z)\nu)-\alpha_a|
       &\le C M_a\delta^{-4}|z|.
\end{align*}
The cubic remainder accumulated over $r+s$ collisions is at most
$C\delta^{-6}|v|^3/\sqrt n$. Since the quadratic form is nonnegative,
split the principal comparison as
$\lambda(z)^{r+s}(A(z)-\alpha_a)
 +\alpha_a(\lambda(z)^{r+s}-e^{-(r+s)z^{\mathsf T}\Gamma_{R,\delta}z/2})$,
where $A(z)$ is its principal amplitude. The first term uses the
amplitude bound above; the second uses $|\alpha_a|\le M_a$ and
$|e^w-1|\le |w|e^{|w|}$. Thus the cubic error is multiplied only
by $C M_a$, not by the multiplier norm. This proves the required spectral bound in the
case of two long blocks.

If $r<q$, delete the first $r$ terms from the real exponent on the left
of \eqref{eq:marked-chronological-pairing}. The cost is at most
$C M_a|v|q/\sqrt n$, by boundedness of $h_{R,\delta}$. The remaining
pairing is $\ell(L_z^s(a_\delta\nu))$. Its principal amplitude differs
from $\alpha_a$ by at most $C M_a\delta^{-4}|z|$, by the smooth-density
embedding, and its complementary term is bounded by
$C M_a\delta^{-2}\rho^s$. Here $s\ge q$. If $s<q$, the remaining
pairing is instead $\ell(M_{a_\delta}L_z^r\nu)$, with exactly the same
bounds. At most one of $r,s$ is less than $q$. Restoring the removed
Gaussian time coefficient changes its value by at most
$C M_a|v|^2q/n$. This handles zero blocks as well; no spectral expansion
of an identity operator is used.

Finally replace $\Gamma_{R,\delta}$ by $\Gamma_R$. The segment joining
these matrices is positive semidefinite, and
\eqref{eq:collision-covariance-limit} bounds the change by
$C M_a|v|^2\delta\ell_\delta$. All the displayed errors combine to
\eqref{eq:marked-spectral-bound}. Only forward collision operators
occurred in the proof.
\end{proof}

\subsection{The growing band and norm growth}
\begin{theorem}[Integrated central band with an arbitrary return mark]
\label{thm:marked-return-band}
There is a constant $C$ such that for every $n\ge2$, $R\in I$,
$0\le k\le n$, and complex $a$ as in \eqref{eq:marked-norms},
\begin{align}
 &\int_{|v|\le2n^{1/200}}
   \left|C^{[k],a}_{n,R}(v)-\alpha_a e^{-v^{\mathsf T}D_Rv/2}\right|\dd v
      \notag\\
 &\hspace{15mm}\le C\left\{
 M_an^{-3/280}\sqrt{\log(2+n)}+V_an^{-9/175}\right\}.
      \label{eq:marked-return-band}
\end{align}
In particular the error tends to zero uniformly for bounded $M_a$ and
$V_a=O(n^\kappa)$ with any $\kappa<9/175$.
\end{theorem}
\begin{proof}
Take $r=m_k$ and $s=m_{n-k}$, so
\[
 (r+s)/n=c_*^{-1}+O(n^{-1}),
\]
with an error independent of $k$. They satisfy the length bounds in
Lemma~\ref{lem:marked-spectral}. Choose
\begin{equation}\label{eq:marked-scales}
 U_n=n^{1/200},\qquad \delta_n=\tfrac14n^{-1/14},\qquad
 q_n=\left\lceil\left(1+\frac4{|\log\rho|}\right)\log(2+n)\right\rceil.
\end{equation}
Then $\rho^{q_n}\le(2+n)^{-4}$ and $2q_n\le r+s$ for sufficiently
large $n$, uniformly in the mark. On $|v|\le2U_n$, both analytic
restrictions hold for large $n$: their positive margins are
$1/2-2/14-1/200$ and $1/2-6/14-3/200$.

Apply the marked spectral estimate and
Proposition~\ref{prop:marked-stopping}. Replacing $(r+s)\Gamma_R/n$
by $D_R$ costs at most $C M_a|v|^2/n$. Integration uses
$\int_{|v|\le2U}|v|^j\dd v=C_jU^{4+j}$. The contribution involving
$V_a$ is only
\[
 C V_a\delta_nU_n^4=C'V_an^{-9/175}.
\]
For clarity, the negative-power margins of the remaining polynomial
terms are displayed below; logarithmic factors are recorded separately.
\begin{center}
\begin{tabular}{lll}
Source & Margin & Logarithmic factor\\ \hline
Observable smoothing & $3/280$ & $\sqrt{\log(2+n)}$\\
Covariance smoothing & $29/700$ & $\log(2+n)$\\
Principal amplitude & $53/280$ & $1$\\
Cubic remainder & $51/1400$ & $1$\\
Clock exception & $9/50$ & $1$\\
Clock window & $7/40$ & $\log(2+n)$\\
Short collision block & $19/40$ & $\log(2+n)$\\
Time coefficient and rounding & $97/100$ & $\log(2+n)$
\end{tabular}
\end{center}
The complementary-power contribution is at most
$C M_an^{2/14+4/200}(2+n)^{-4}$ and is smaller than the displayed
polynomial errors. The first margin in the table is strictly smaller
than all the others, so it absorbs their logarithmic factors.
This proves \eqref{eq:marked-return-band} for large $n$. For the
remaining finite range, both transforms have modulus at most $M_a$,
and the frequency ball has finite volume; enlargement of $C$ proves
the stated estimate for all $n\ge2$. The last assertion follows by
substitution, without changing the smoothing scale with the insertion.
\end{proof}

\begin{corollary}[Conditioning on a state at a marked return]
\label{cor:marked-state-conditioning}
Let $E_{n,R}\subset Y_R^*$ have positive measure and a zero-extended
indicator in $\mathrm{BV}(M)$. Set
$V_{n,R}=\|\one_{E_{n,R}}\|_{\mathrm{BV}}$. Under the section probability,
condition on the exact event
\[
 A_{n,k,R}=\{x:(F_R^*)^kx\in E_{n,R}\}\n\]
Uniformly for $0\le k\le n$,
\begin{align}
 &\int_{|v|\le2n^{1/200}}
 \left|\mathbb E_{\nu_R^*}\left[
 e^{iv\cdot(J_{n,R}-n\bar G_R)/\sqrt n}\mid A_{n,k,R}\right]
          -e^{-v^{\mathsf T}D_Rv/2}\right|\dd v\notag\\
 &\quad\le\frac{C}{\nu(E_{n,R})}
  \left(n^{-3/280}\sqrt{\log(2+n)}+V_{n,R}n^{-9/175}\right).
       \label{eq:marked-conditional-band}
\end{align}
For example this tends to zero if, uniformly in $R$,
$\nu_R^*(E_{n,R})\ge c n^{-\beta}$ and $V_{n,R}\le Cn^\kappa$,
where $\beta<3/280$ and $\beta+\kappa<9/175$.
\end{corollary}
\begin{proof}
Invariance gives $\nu_R^*(A_{n,k,R})=\nu(E_{n,R})/c_*$.
The conditional characteristic function is exactly
$C^{[k],\one_{E_{n,R}}}_{n,R}/\nu(E_{n,R})$. Apply
Theorem~\ref{thm:marked-return-band} with $M_a=1$. The two strict
exponent inequalities absorb the logarithm. This argument uses the
same event throughout and makes no comparison with another event.
\end{proof}

\paragraph{Relation to the raw density theorem.}
These estimates control the central integral of the actual complex
measures \eqref{eq:marked-measure}. The exact local decomposition
\eqref{eq:exact-local-decomposition} also applies to these measures by
linearity. More explicitly, suppose $D_R\ge d_0 I_4$; suppose they admit
extracted measures $E^{[k],a}_{n,R}$ with mixed densities; and suppose
the residual transform $H^{[k],a}_{n,R}$ is in
$L^1(\T^3\times\R)$, satisfies the complementary-tail condition
\eqref{eq:new-complementary-tail}, and satisfies the local edge
conditions \eqref{eq:local-edge-errors}, all uniformly in the chosen
marks and insertions. If the right-hand side of
\eqref{eq:marked-return-band} tends to zero uniformly, the proof of
Corollary~\ref{cor:raw-from-growing-arc} gives its raw density conclusion
for these marked measures. Thus its central hypothesis is now proved
also for terminal and single intermediate state insertions. The other
hypotheses have not been proved by the recentering argument.

The band still has physical radius $n^{-99/200}$, not $n^{-2/5}$.
Multiple separated return-state factors are not one insertion at a
single mark, and composing them with a long induced iterate does not
give the variation norm used here. Likewise an exact final lattice or
flight-time constraint is a condition on the record, not generally a
collision-state set $E_{n,R}$ of the class in
Corollary~\ref{cor:marked-state-conditioning}. Its relative error bound
is consequently not a replacement for the weighted raw local limits
needed for those constraints. Nor does a two-sided orbit identity
provide a pointwise representative of the zero-variance coboundary.
The covariance nondegeneracy, complete complementary integral, and
full critical/singular residual estimates retain their separate roles.
